% Generated draft: requires mathematical and source review.
\documentclass[11pt]{article}
\usepackage[margin=25mm]{geometry}
\usepackage{fontspec}
\setmainfont{Latin Modern Roman}
\usepackage{amsmath,amssymb,mathtools,unicode-math}
\setmathfont{Latin Modern Math}
\usepackage{enumitem,needspace,xurl,hyperref}
\setlength{\parindent}{0pt}
\setlength{\parskip}{5pt}
\setlength{\emergencystretch}{3em}
\newcommand{\N}{\mathbb{N}}
\newcommand{\Z}{\mathbb{Z}}
\newcommand{\Q}{\mathbb{Q}}
\newcommand{\R}{\mathbb{R}}
\newcommand{\C}{\mathbb{C}}
\newcommand{\F}{\mathbb{F}}
\newcommand{\A}{\mathbb{A}}
\newcommand{\PP}{\mathbb{P}}
\newcommand{\E}{\mathbb{E}}
\newcommand{\eps}{\varepsilon}
\newcommand{\dd}{\,\mathrm{d}}
\begin{document}
\section*{UnsolvedMath ID: 10300045}
\textbf{Source problem code:} catalogue entry
\subsection{Definitions and mathematical statement}
Let \(\mathcal F\) be a taut codimension-one foliation of a hyperbolic \(3\)-manifold \(M\). Under the source's asserted identification \(\widetilde M\cong\mathbb H^3\), let \(\widetilde{\mathcal F}\) be the lifted foliation. Its leaf space is the quotient
\[
 L=\widetilde M/{\sim},\qquad x\sim y\ \Longleftrightarrow\ x,y\text{ lie on the same leaf of }\widetilde{\mathcal F},
\]
with quotient map \(\pi:\widetilde M\to L\). The foliation is \(\R\)-covered when \(L\) is homeomorphic to \(\R\). Leaf spaces here can be non-Hausdorff; ordinary end and recurrence definitions therefore cannot be imposed without checking the source conventions.

Question 10.7 asks the following.
\begin{enumerate}
\item For a random walk \(\gamma\) in \(\widetilde M\), describe the typical behavior of its projected trajectory \(\pi(\gamma)\).
\item Does such a random walk converge to a definite end of \(L\)?
\item Does the analogous conclusion hold for a random geodesic?
\item Is the asymptotic projected behavior well defined under bounded-distance perturbations? More precisely, is there a full-measure set \(A\) of geodesics such that, for every \(\gamma\in A\) and every admissible \(\gamma'\) at bounded distance from \(\gamma\), the trajectories \(\pi(\gamma)\) and \(\pi(\gamma')\) exit the same end of \(L\)?
\item If \(\mathcal F\) has one-sided branching, does a random walk always exit in the unbranching direction?
\item If there is positive probability that a random walk exits a proper positive or negative end of \(L\), must \(\mathcal F\) be \(\R\)-covered? What analogous conclusion holds when there is positive probability of recurrence?
\end{enumerate}
These questions require a law for the random walk, a measure for random geodesics, definitions of ends and recurrence in \(L\), and the source's meanings of one-sided branching and proper positive/negative ends. In the fourth question, the quantifier over all bounded-distance perturbations is inside the full-measure geodesic condition; it is not merely an almost-everywhere statement about pairs. The meaning of bounded distance and the class of perturbations are unspecified. The word ``always'' in the fifth question also requires an almost-sure versus pathwise interpretation.
\subsection{Short English statement}
Project random walks and random geodesics in the universal cover of a tautly foliated hyperbolic three-manifold to its leaf space. Determine their limiting ends, stability under bounded-distance perturbations, and what branching, positive escape probabilities, or recurrence force about the foliation. The supplied record leaves the stochastic laws and several leaf-space conventions implicit.
\subsection{Sources}
\begin{itemize}
\item Local UnsolvedMath registry, permanent ID \(10300045\), exact source code \texttt{catalogue entry}, ``Hyperbolic geometry --- Question 10.7.'' No individual proposer or proposed year is supplied.
\item Danny Calegari, \emph{Problems in foliations and laminations of 3-manifolds} (2002), Question 10.7: \url{https://arxiv.org/abs/math/0209081}. Author-source list containing 67 numbered questions.
\item Source-TeX extraction, accessed 2026-07-29. Source status NEEDS\_REVIEW: presented as open, with a current primary-literature audit still required. The supplied OPEN-TRIAGE summary mentions random walks, boundary actions, and Poisson boundaries of foliations, but provides no additional specific references and reports no verified answer to this question.
\item Dataset attribution: Ulam AI / UnsolvedMath Contributors, CC BY 4.0. Source rights NEEDS\_REVIEW: publicly accessible author/publisher copy, redistribution terms requiring review.
\end{itemize}
\end{document}
